\section{The congruence ideal at a type prime}

% BEGIN P84D1
\paragraph{P84D1. The ideal on the retained points.}
Write $O$ for the coefficient discrete valuation ring, $p$ for
the residue characteristic, $e=\ell_O(O/pO)$, and $X$ for
the original finite set of Hecke points at the chosen index.
On a good cut let $S\subset X$ be the retained set and put
$Y=S\cap Z=Z_{\rm good}$, $B=S\setminus Y$.
For $U\subset X$ let $T_U$ denote the image of the original
Hecke algebra $T$ in $O^U$. For every $\pi\in Y$ put
\[
 J_\pi=\{h(\pi):h\in T,\ h(b)=0\text{ for every }b\in B\},
 \qquad c_\pi=\ell_O(O/J_\pi)/e .
\]
This set of values is an ideal since $T$ contains the scalars.
Restriction onto $T_S$ is surjective, and vanishing off $Y$ in
$S$ means exactly vanishing on $B$. Thus $J_\pi$ equals the
native evaluation ideal on the cut. In $X$ the complement of $B$
is $Z\cup(X\setminus S)$, the native cut locus, so the same
$c_\pi$ also equals the original $\mathrm{vccCut}(i,S,\pi)$.

Apply the coordinate filtration inequality LN4 to $Y$ on the
cut lattice. P84C4 identifies the good and theta subsets there.
Equivalently the general endomorphism index bound applies with
the identity and scalar $1$, whose valuation is zero. With $m_\pi$
the cut eigenspace multiplicity this gives
\[
                  \eta'_{\rm loc}(i)\le\sum_{\pi\in Y}m_\pi c_\pi .
\]
This use of the identity is in the general length inequality.
P84C4 separately supplies the prescribed group-cardinality
scalar for the commuting condition in the family record.
% END P84D1

% BEGIN P84D2
\paragraph{P84D2. An exact calculation from two integral restriction images.}
Write $d(U)=\ell_O(O^U/T_U)$, with $d(\varnothing)=0$.
The separation condition on $T$ makes every such length finite.
Indeed for each $x\in U$ and each $y\ne x$ in $U$ choose a
Hecke function taking distinct values at $x,y$. Subtract its
value at $y$ and multiply these functions over $y$. The result
has support $x$ on $U$ and nonzero value at $x$. For
$U=\{x\}$ the empty product is $1$. These elements span the
coordinate lines over the fraction field, and $O$ is Noetherian.
For every $\pi\in Y$ we have
\[
           e c_\pi=d(B\cup\{\pi\})-d(B).
\]
In fact the map $O\to O^{B\cup\{\pi\}}/T_{B\cup\{\pi\}}$
supported at $\pi$ has kernel $J_\pi$. Its image is the kernel
of restriction to $O^B/T_B$: a representative in that kernel
agrees on $B$ with some original Hecke function, whose subtraction
leaves just its coordinate at $\pi$. Restriction is surjective
by extension of functions. Additivity of lengths proves the
formula. Each $d(U)$ is the sum of the valuations of the
Smith invariant factors of a matrix of integral module generators
of $T_U$, with the uniformizer of $O$ assigned value $1$.

There is also a bound needing only pairs. Let $a_{\pi b}$ be
the least valuation of a nonzero $h(\pi)-h(b)$ for $h\in T$.
It exists by separation and discreteness. Multiplying one
minimizing difference function $h-h(b)$ for each $b\in B$
produces a function whose value at $\pi$ has valuation
$\sum_{b\in B}a_{\pi b}$. Therefore
$e c_\pi\le\sum_{b\in B}a_{\pi b}$. In particular the
congruence lengths vanish if $B$ is empty, or, more generally,
whenever each $b\in B$ can be distinguished from $\pi$
modulo the maximal ideal. On a general good cut the set $B$
may contain points, and the exact difference above tests their
contribution.
% END P84D2

% BEGIN P84D3
\paragraph{P84D3. The point bound implies the paired condition with the same constant.}
Consider any of the tests in
$\mathrm{PairedCharacterConditionAlt}(C,C_T)$ on a good cut.
In particular the two type primes are distinct, $p=g_k$, the
size, sign and host-jump hypotheses hold, and the specified
rank Selmer and cyclotomic lengths are finite. Write $n=r+t$,
$f=[E^+:\mathbb Q]$. Denote by $A_\pi$ the part of the
native right hand side independent of the table constant and
of the quadratic reserve:
\[
 A_\pi=s_\pi+\sum_{0<j<t}\mathrm{shaLen}_p(j)
        +v_p(h_E)+4n v_\delta+t\sum_{q\in Q_i}v_\zeta(q).
\]
Suppose that at each $\pi\in Y$ the length from P84D1 satisfies
the native point bound, namely there are integers $2\le N\le n+1$,
$1\le u$, $4u\le N^2$, for which
\[
 c_\pi\le 4 f n C_T\bigl(Nu+(\eta+t+1)(N+u)\bigr)+A_\pi .
\]
The witnesses may vary with $\pi$. For every $C_T\in\mathbb N$
this implies the paired inequality with $C=C_T$.
Indeed the native rank assumptions imply $n\ge3$. Then
$N\le4n/3$, $u\le4n^2/9$, $N+u\le8n^2/9$ and
$Nu\le16n^3/27$. Since $N+u\le 8n^3/27$ as well, we have
\[
 Nu+(\eta+t+1)(N+u)
 \le \tfrac{24}{27}n^3+\tfrac89 n^2(\eta+t)
 \le n^2(n+\eta+t).
\]
The native right hand side adds $2fr^2$ to
$4fnC_Tn^2(n+\eta+t)+A_\pi$. Multiply by the
nonnegative $m_\pi$, sum, and use P84D1. This gives the
precise conclusion for the tested families and parameters.
If the point bound holds for every native test, the paired
condition follows with the same uniform constant. Neither
the argument nor the alternate test requires an extra
$\mathrm{SetSQ}$ assumption. P84D2 also proves the same
conclusion directly whenever its Smith differences divided by
$e$ are bounded by the native right hand side at every test.
% END P84D3

% BEGIN P84D4
\paragraph{P84D4. A cut carrying one prescribed tuple of local types.}
Work on the good cut of P84B1, writing $X$ for its spectrum
and $Y=Z=Z_{\rm good}$. Fix a type-prime index and $\pi\in Y$.
Its good list supplies a parameter $\sigma$ whose polynomial
labels form one tuple $a$. On every representative above $\pi$,
the same-list comparison of P84C2 puts its whole type in
$\mathrm{MtPar}(\sigma)$. By the proof of P84C2, $E(a)$ fixes
every such vector. Use just the one integral central element
lifting $N_G E(a)$ from P84C1 and let $\beta$ be its action
on the cut lattice. It commutes with Hecke. The transported
$\mathrm{EigSpan}$ comparison gives $\beta=N_G$ on the entire
rational eigenspace at $\pi$.

Take $S_\pi$ to consist of those $x$ admitting a representative
and a type vector on which $E(a)$ acts nontrivially. This
contains $\pi$, since a representative with nonzero type exists.
At every $x\notin S_\pi$ all representatives have zero action,
so $\beta$ vanishes on the whole rational eigenspace there.
Also $S_\pi\cap Z\subseteq Z_{\rm good}$ since this is a good
cut. The integer $b=N_G$ is nonzero and
$\ell_O(O/bO)/e=v_p(N_G)$ by factoring it in the coefficient
DVR. These are precisely the conditions of
$\mathrm{CutAt}(i,k,\pi,S_\pi,\beta,b)$.

For $x\in S_\pi$ retain the same witnessing representative
and vector. Each of the commuting local series factors of
$E(a)$ acts nontrivially on that vector. Expanding a local
series into its finite sum of irreducible character projectors
shows that some one of these projectors acts nontrivially.
The local type occurrence identity in the given depth-zero
record then puts that irreducible in the representative's
specified set of types. Its label is the prescribed polynomial
at that place. At places over $g_k$ these polynomials are
squarefree and avoid both $\chi_J$ and its inverse by P84C3.
Thus the single tuple of polynomials, followed at each $x$ by
one representative before all places, witnesses the full
$\mathrm{TypeLocus}(P,i,k,S_\pi)$. The set $S_\pi$ is chosen
before the good Galois family or Selmer data tested at $\pi$.
% END P84D4

% BEGIN P84D5
\paragraph{P84D5. The cost of choosing a different local type at each point.}
For each $\pi\in Y$ use P84D4's cut or any native
$\mathrm{CutAt}(i,k,\pi,S_\pi,\beta_\pi,b_\pi)$. Put
$B_\pi=S_\pi\setminus Z$. Set $c_\pi^*=
[d(B_\pi\cup\{\pi\})-d(B_\pi)]/e$ from the algebra on
this spectrum. P84D2 identifies it with
$\mathrm{vccCut}(i,S_\pi,\pi)$, since only vanishing on
$B_\pi$ is tested. Then
\[
 \eta_{\rm loc}\le\sum_{\pi\in Y}m_\pi
            \bigl(c_\pi^*+v_p(N_G)\bigr).
\]
To see this even when all the supports are different, use LN4's
filtration of the projected lattice $U=P_Y M$ modulo
$M\cap V_Y$, where $V_Y$ is the sum of the rational eigenspaces
in $Y$. At its step for $\pi$, choose $h\in T$ vanishing on
$B_\pi$ whose value at $\pi$ generates the evaluation ideal.
Existence follows from P84D1 with this vanishing set. The
endomorphism $h\beta_\pi$ is integral, kills every rational
eigenspace outside $Y$, and acts at $\pi$ by $h(\pi)b_\pi$.
Indeed $\beta_\pi$ preserves each eigenspace by commutation,
and outside $Y$ either it vanishes or $h$ does. Therefore
$h\beta_\pi$ sends $U$ into $M\cap V_Y$: for $P_Ym\in U$ its
image equals $h\beta_\pi m$. It preserves the partial sums of
eigenspaces in the filtration. On the current graded piece
it acts by $h(\pi)b_\pi$, which consequently kills the
graded quotient. That quotient is generated by at most $m_\pi$
elements by the free-lattice argument of LN4. Summing its
length bounds and using the scalar length test in
$\mathrm{CutAt}$ proves the inequality.

At a paired test let $g=p$ and $g'$ be the other generator
prime. The logarithm guard gives $g'\le g^2$. There are $f$
factors at each prime in $N_G$. The elementary column count for
$\mathrm{GL}_r(\mathbb F_q)$ gives its order as
$q^{r(r-1)/2}\prod_{a=1}^r(q^a-1)$. For the $g$ factors
only the initial power contributes. For $g'\ne g$, each
positive integer $g'^a-1<g^{2a}$ has $g$-valuation at most
$2a$. Hence
\[
 v_g(N_G)\le f\{r(r-1)/2+r(r+1)\}\le 2fr^2.
\]
This is the quadratic reserve in P84D3. Apply its polynomial
estimate to $c_\pi^*$ and absorb this scalar cost before summing
multiplicities. The native point bounds on these smaller
$S_\pi$ now suffice for $\mathrm{PairedCharacterConditionAlt}$
with $C=C_T$. In particular P84B2 applies with these
pointwise tests in place of the tests on the whole retained
spectrum. The Smith differences for the larger bad set need
not satisfy the smaller-set bounds.
% END P84D5
